This report documents the development, calibration, and business evaluation of an
investment recommendation system for FinPUC, a fictitious financial advisory
platform that generates personalized portfolios for clients in the US stock
market. The system serves five risk profiles and pursues a single objective:
maximizing the client's expected return. The firm's revenue is measured
afterwards, as a consequence of the active managed balance, through a fixed
monthly fee of 0.5\% on that balance, with no transaction or rebalancing fees.

The investment universe comprises 601 stocks filtered from 1,406 tickers. Three
methodological layers were implemented: Markowitz mean-variance optimization as
the base case, four Black-Litterman view families calibrated independently, and a
Monte Carlo layer that simulates client behavior. Each view is tuned over a
three-horizon sequential protocol (calibration, validation, and clean test) to
avoid contaminating the evaluation period. Client behavior is modeled with a
single weekly abandonment probability; recommendations are issued semiannually
while abandonment is evaluated weekly.

Among the calibrated views, general momentum delivers the largest economic gain:
it improves the P4 score by $+12.4\%$ under the robust median aggregation and by
$+85.7\%$ under the average aggregation, while also improving Sharpe by
$+5.2\%$ relative to the Markowitz base. This advantage comes with higher
turnover (34.6\%) and sector concentration (HHI 0.261), while the
unemployment-macro view provides superior operational stability (7.3\% turnover,
HHI 0.157). The projected-return panorama analysis confirms that the
recommendation does not depend on a single expected-return point. The
recommended methodology is Black-Litterman with the general-momentum view for
aggressive profiles, with the unemployment-macro view as a defensive alternative
for conservative profiles.

\textbf{Keywords:} portfolio optimization, Markowitz, Black-Litterman
calibration, Monte Carlo simulation, recommender system, weekly abandonment,
managed-balance fee, rolling validation, CVXPY.
